% One-time imported draft. Editable native source; no Markdown regeneration.
\section{Assign meanings to registers at
  boundaries}\label{007-bresenham-line-p03-assign-meanings-to-registers-at-boundaries}

The working point changes as the loop progresses, while the ordered
endpoint stays fixed. The horizontal and vertical differences also stay
fixed. Keeping these meanings distinct makes the branch tests readable:
a coordinate comparison determines whether the line is complete; an
error comparison determines which coordinate changes next.

The signed y direction is separate from its magnitude. A difference of
negative three means a distance of three rows with a step of negative
one. This separation permits the same threshold logic to serve rising
and falling lines in screen coordinates.

The table of registers describes the reference program, not a required
student ABI. A free-writing solution may allocate them differently. The
supplied input block and framebuffer layout are the interface that other
tools and grading scripts use.
